% Ad Familiares 7.26 --- headnote
% ad-familiares-07-26, 46 BC (intercalary month), Tusculum

\textit{Cicero to M. Fadius Gallus, written from the Tusculan villa
in the intercalary month following 46 BC, and one of the most
quoted of all his letters. Cicero has fled to Tusculum after ten
days of bowel-trouble in town, where he could not persuade his
clients he was really unwell because he had no fever; two days of
strict fasting --- not even water tasted --- has left him spent.
The body of the letter is a model of the comic juristic-medical
register: Cicero dreads, he says, every illness, but particularly
that on which the Stoics fasten Epicurus, who admitted that
``strangurical and dysenteric afflictions'' troubled him --- the
one a complaint of greed, the other of a more disgraceful
intemperance still. Three of the four Greek phrases are clinical
terms transliterated straight: \textit{strangourika kai dysenterika
pathē}, \textit{dysenterian}, \textit{diarroia}; the fourth,
\textit{litotēta}, ``frugality,'' is the elegantly philosophical
label that the joke pivots on.}

\textit{The joke: the sumptuary law (Caesar's of 46) capped
expenditure on meat and fish but, in good legalese, exempted
\textit{terra nata}, ``things born of the earth.'' Rome's
fashionable gourmets accordingly competed to season mushrooms,
herbs, wild greens, and roots into delicacies more delicious than
anything banned. At an augural dinner at Lentulus's, Cicero ---
who could pass up oysters and lampreys without difficulty --- was
ambushed by beet and mallow; the result was as advertised. The
Perseus dateline is \textit{in m. intercalari post a. 708 (46)},
the intercalary month after the regular year; the meta entry's
\textit{-0046-07-27} placeholder is a year-month-precision stub.
Anicius, who had called and ``seen me being sick,'' had carried
the news to Gallus; the closing is half a thank-you note, half
the resolve, decidedly belated, to be more careful.}
